% TOP_PROBLEM: {"id": 3317, "title": "3-Colourability of Arrangements of Great Circles", "category_id": 3, "category": {"id": 3, "name": "graph_theory", "display_name": "Graph Theory", "description": "Problems involving graphs, networks, and their properties.", "slug": "graph-theory", "order_index": 3, "created_at": "2026-07-31T15:26:25.671Z"}, "status": "open", "problem_number": "OPG-34915", "set_id": 12}
% FIRST_POSTED: 2026-10-02
% ATTEMPT_STATUS: unresolved
% MODEL: GPT 6 Astra Ultra
% UnsolvedMath ID 3317; source catalogue code OPG-34915
% !TeX program = lualatex
% Research attempt dated 2026-10-02; canonical statement preserved.
\documentclass[11pt,a4paper]{article}
\usepackage[margin=25mm]{geometry}
\usepackage{fontspec}
\setmainfont{lmroman10-regular.otf}[BoldFont=lmroman10-bold.otf,
 ItalicFont=lmroman10-italic.otf,BoldItalicFont=lmroman10-bolditalic.otf]
\usepackage{amsmath,amssymb,mathtools}
\usepackage{unicode-math}
\setmathfont{latinmodern-math.otf}
\AtBeginDocument{\let\setminus\smallsetminus}
\usepackage{xurl,hyperref}
\hypersetup{colorlinks=true,urlcolor=blue}
\setcounter{secnumdepth}{0}
\setlength{\parindent}{0pt}
\setlength{\parskip}{5pt}
\setlength{\emergencystretch}{3em}
\begin{document}
\section*{3-Colourability of Arrangements of Great Circles}\label{3317}
{\small UnsolvedMath ID 3317 · OPG-34915 · Graph Theory · OpenGarden}
\subsection{Definitions and mathematical statement}
Let $S^2=\{x\in\mathbb R^3:\|x\|=1\}$. A great circle is the
intersection of $S^2$ with a two-dimensional linear subspace of
$\mathbb R^3$. Fix a finite collection $\mathcal C$ of $m\ge2$
distinct great circles in general position in the sense that no
point of the sphere belongs to three of them. Each pair meets at
two antipodal points.

The arrangement graph has these intersection points as vertices.
For each circle, join two consecutive intersection points along
that circle by the intervening circular arc, provided that arc has
no other intersection point in its interior. These arcs are the
edges of the embedded graph. The graph is 4-regular when parallel
edges are counted; in the two-circle case parallel edges do occur
and impose the same colouring constraint as a single edge.

The conjecture asks whether this arrangement graph always has a
proper 3-colouring: a map from its vertices to $\{1,2,3\}$ taking
different values at the endpoints of every arrangement edge.
No relation is required between colours at antipodal vertices.
For fewer than two circles there are no intersection vertices,
so the colouring requirement is vacuous.

The restriction to great circles is a geometric hypothesis, not
merely that the graph is planar and 4-regular. Likewise, the
no-triple-intersection assumption specifies which arrangement graph
is intended. The target is ordinary vertex colouring from one
three-colour palette; the stronger list-colouring question mentioned
in the source is separate.
\subsection{Short English statement}
Is the intersection graph drawn along any general-position arrangement of great circles properly vertex-colourable with three colours?
\subsection{Research attempt}
\paragraph{Scope and method.}
This research attempt by GPT-6 Astra Ultra is dated 2 October 2026.
It preserves the catalogue statement and distinguishes elementary proofs
below from cited prior work. No novelty claim or formal proof-assistant
verification is made. The final paragraph states the precise remaining scope.
The finite checks described below are reproducible in
\texttt{scripts/check\_attempt\_batch03\_root\_2026\_10\_02.py}.
\paragraph{A small-arrangement strategy.}
Give each intersection the label of the unordered pair of circles
that contains it. There are two antipodal vertices of each label.
A sufficient condition for a colouring is to give the edges of the
complete graph on the circle labels a proper edge-colouring, and
transfer each pair's edge colour to both of its intersection points.
This condition deliberately imposes more than the original problem:
the original allows different colours at antipodes and only tests
consecutive intersections along each circle.

\paragraph{Why the transfer works.}
With at least three circles, two consecutive intersections on a
circle cannot both have the same pair label. Those two intersections
would be antipodal. Every third great circle meets the interior of
each of the two open semicircles between them, contradicting
consecutiveness. Thus adjacent vertices along circle $i$ have labels
$\{i,j\}$ and $\{i,k\}$ with $j\ne k$. A proper edge-colouring
of the complete label graph gives these different colours. This
argument uses the antipodal intersection property, rather than
planarity alone.

\paragraph{Exact answers through four circles.}
Two circles give two vertices joined by four parallel arcs, so two
colours suffice and one does not. With three circles, give the
three pair labels three different colours. The resulting graph is
the octahedral graph $K_{2,2,2}$ and is properly 3-coloured.
With four circles, assign colour A to pairs $12,34$, colour B to
$13,24$, and colour C to $14,23$. Each circle occurs in one pair
of each colour. The transfer argument proves a proper colouring
for every general-position arrangement of four circles, without
needing to classify their possible geometries.

There must in fact be a triangle in every arrangement with at least
three circles. Write $m$ for the number of circles. The connected
spherical graph has $V=m(m-1)$ vertices and $E=2V$ edges. It is
simple: two different circles can share only their antipodal
intersection vertices, and another circle interrupts both arcs
between such vertices; along a single circle there are at least
four distinct intersection vertices. Every edge belongs to a
circle cycle, so there is no bridge. Euler's formula gives $F=V+2$.
If every face had length at least four, counting edge sides would
give $2E\ge4F$, or $4V\ge4V+8$, a contradiction.
Hence some face has length three. In particular the three- and
four-circle answers are exactly three, rather than just upper bounds.

\paragraph{The obstruction in the label method.}
For $m\ge5$, a proper edge-colouring of the complete label graph
needs at least $m-1\ge4$ colours at each label vertex. This only
shows the failure of our sufficient condition. It does not show
that the arrangement needs four colours: many pairs incident with
a circle correspond to nonconsecutive vertices and need not be
different in a valid arrangement colouring. A general argument
must exploit that cyclic order, or allow antipodes different colours,
without creating a conflicting constraint on another circle.

\paragraph{Literature caution and verification.}
Cahit's 2004 preprint [R1] announces a full proof. The discussion in
[S2] explicitly disputes that proof, and the later problem collection
[R2, Problem 44] retains the question. We therefore do not treat the
announcement as a verified solution. The root script checks the
pair-label construction against every possible local ordering of
the other circle labels for $m=3,4$. These tests even allow orders
that may not be simultaneously geometrically realizable; validity
on that larger family supports the sufficient-condition lemma but
does not extend its number of colours.
The cases of at most four circles are established here. No general
three-colouring theorem is proved by this attempt.

\subsection{Sources}
\begin{itemize}
\item[S1] ULAM.AI and the UnsolvedMath Contributors, supplied problems.json, record 3317 (OPG-34915), OpenGarden collection. Catalogue identity and source transcription; CC BY 4.0.
\url{https://huggingface.co/datasets/ulamai/UnsolvedMath}.
\item[S2] Open Problem Garden, 3-Colourability of Arrangements of Great Circles. Original problem and discussion; the mathematical formulation below is expanded from the supplied source record.
\url{http://www.openproblemgarden.org/op/3_colourability_of_arrangements_of_great_circles}.

\item[R1] I. Cahit, Three Colorability of an Arrangement Graph of Great
Circles, arXiv:math/0408363. Claimed proof;
the dispute in [S2] is material to its use here.
\url{https://arxiv.org/abs/math/0408363}.
\item[R2] The Open Problems Project, Problem 44, Great-Circle Colorings.
\url{https://topp.openproblem.net/master.pdf}.

\end{itemize}
\end{document}
